\section*{Proof Development Record}

This note records the development of the index-free precursor separately from the formal manuscript. Symbols follow \texttt{preamble.tex}. The missing repair-order and recursion details were subsequently resolved in \texttt{notes/formal\_proof\_audit.md}. The manuscript now states the strictly sharper FIFO-indexed theorem; its additional derivation is recorded in \texttt{notes/indexed\_bound\_derivation.md}.

\subsection*{Local repair construction}

Let \(a\) and \(a'\) be consecutive vehicles from the same approach and the same proposed platoon. Consider a FIFO-feasible sequence
\begin{equation}
s=[P,a,B,a',R],
\end{equation}
where \(B=(b_1,\ldots,b_m)\) is nonempty. Because \(a'\) is the immediate FIFO successor of \(a\), no vehicle in \(B\) belongs to their approach. Construct
\begin{equation}
s'=[P,a,a',B,R].
\end{equation}
Let \(T=C_a(s)=C_a(s')\), \(d=r_{a'}-r_a\), \(q=\max\{d,\hF\}\), \(r=|R|\), and \(M=m+r+1\).

\subsection*{Passing time of the moved vehicle}

Since \(T\ge r_a\),
\begin{equation}
r_{a'}=r_a+d\le T+d.
\end{equation}
Therefore,
\begin{equation}
C_{a'}(s')=\max\{r_{a'},T+\hF\}\le T+q.
\end{equation}
In the original sequence,
\begin{equation}
C_{b_m}(s)\ge T+\hS+(m-1)\hF
\end{equation}
because the transition from \(a\) to \(b_1\) requires \(\hS\), and every subsequent transition requires at least \(\hF\). Since \(b_m\) and \(a'\) belong to different approaches,
\begin{equation}
C_{a'}(s)\ge C_{b_m}(s)+\hS
\ge T+2\hS+(m-1)\hF.
\end{equation}
Hence
\begin{equation}
C_{a'}(s')-C_{a'}(s)
\le q-2\hS-(m-1)\hF.
\label{eq:note-moved-change}
\end{equation}

\subsection*{Passing times of the intermediate block}

For \(b_1\),
\begin{align}
C_{b_1}(s')
&=\max\{r_{b_1},C_{a'}(s')+\hS\}\\
&\le\max\{r_{b_1},T+q+\hS\}\\
&\le C_{b_1}(s)+q.
\end{align}
The order and pairwise separation requirements inside \(B\) are unchanged. By induction,
\begin{equation}
C_{b_j}(s')\le C_{b_j}(s)+q,
\qquad j=1,\ldots,m.
\end{equation}
Therefore,
\begin{equation}
\sum_{j=1}^{m}\left[C_{b_j}(s')-C_{b_j}(s)\right]\le mq.
\label{eq:note-block-change}
\end{equation}

\subsection*{Passing times of the unchanged suffix}

Let \(x\) be the first vehicle in \(R\), when \(R\) is nonempty. The new predecessor of \(x\) is \(b_m\), while its old predecessor is \(a'\). Using \(C_{b_m}(s')\le C_{b_m}(s)+q\),
\begin{align}
&C_{b_m}(s')+h(b_m,x)
-\left[C_{a'}(s)+h(a',x)\right]\\
&\le q-\hS+h(b_m,x)-h(a',x)\\
&\le q-\hF.
\end{align}
The last inequality uses \(h(b_m,x)\le\hS\) and \(h(a',x)\ge\hF\). The suffix order is unchanged, so the positive delay perturbation cannot grow under the maximum recursion. Thus
\begin{equation}
\sum_{v\in R}\left[C_v(s')-C_v(s)\right]
\le r(q-\hF).
\label{eq:note-suffix-change}
\end{equation}

\subsection*{Candidate local bound}

Adding Eqs.~\eqref{eq:note-moved-change}--\eqref{eq:note-suffix-change},
\begin{align}
J(s')-J(s)
&\le q-2\hS-(m-1)\hF+mq+r(q-\hF)\\
&=(m+r+1)(q-\hF)-2(\hS-\hF)\\
&=M\pospart{d-\hF}-2(\hS-\hF).
\end{align}

Using the positive part provides a nonnegative per-repair bound. Repeating the repair across all platoon internal links and applying a telescoping sum yields the sequence-specific bound recorded in the formal proof audit. The historical substitution \(M_e\le N\) gives the index-free bound; the manuscript now uses the sharper substitution \(M_e\le N-i_e\).

\subsection*{Audit disposition and remaining research questions}

\begin{enumerate}
  \item Preservation of previously established platoon adjacencies is proved by processing links from left to right within each platoon.
  \item The suffix inequality is proved using max-recursion monotonicity and the fact that \(q-\hF\ge0\).
  \item Determine whether taking a positive part separately for every repair causes avoidable looseness.
  \item Search for cases in which a different repair direction produces a smaller bound.
  \item The FIFO-indexed zero-loss condition is a valid corollary of the formal theorem.
  \item Construct tight or near-tight instances and identify the limiting inequalities.
\end{enumerate}
